% !TeX root = ../main.tex

\documentclass[../main.tex]{subfiles}
\begin{document}

\section{Tổng quan về đối tượng công nghệ}

\subsection{Giới thiệu chung về hệ thống cấp dầu nhiên liệu khởi động}

Trong các nhà máy nhiệt điện đốt than, dầu nhiên liệu khởi động được sử dụng trong giai đoạn khởi động lò, duy trì ngọn lửa mồi, hỗ trợ bắt cháy khi phụ tải thấp và làm nguồn dự phòng khi hệ thống cấp than chưa sẵn sàng hoặc khi điều kiện vận hành yêu cầu chuyển đổi chế độ cháy. Đối tượng công nghệ của bài tập lớn là hệ thống cấp và điều khiển lưu lượng dầu nhiên liệu khởi động cho lò than phun. Đây là một hệ thống công nghệ có tính liên tục, làm việc trong điều kiện nhiệt độ cao, dầu có độ nhớt lớn và yêu cầu an toàn cháy nổ nghiêm ngặt.

Về bản chất, hệ thống này thực hiện nhiệm vụ tiếp nhận, lưu trữ, gia nhiệt, tuần hoàn, đo lường và phân phối dầu đến các đầu đốt của lò hơi. Trong chu trình đó, lưu lượng dầu cấp là đại lượng có ảnh hưởng trực tiếp đến cường độ cháy, nhiệt lượng sinh ra trong buồng lửa và khả năng duy trì trạng thái cháy ổn định. Do đó, bài toán điều khiển lưu lượng dầu không chỉ là bài toán tự động hóa đơn thuần mà còn là một khâu kỹ thuật quan trọng quyết định độ tin cậy của toàn bộ quá trình khởi động lò.

Hình~\ref{fig:hmi_fuel_oil} là giao diện giám sát điển hình của hệ thống cấp dầu nhiên liệu trên DCS. Hình này cho thấy các cụm bể chứa, đường hồi lưu, bộ gia nhiệt và các nhánh cấp cho đầu đốt được tích hợp trong một cấu trúc công nghệ khép kín.

\begin{figure}[htbp]
    \centering
    \includegraphics[width=0.90\textwidth]{Hinhve/hmi-fuel-oil.jpg}
    \caption{Giao diện giám sát hệ thống cấp dầu nhiên liệu khởi động cho lò hơi}
    \label{fig:hmi_fuel_oil}
\end{figure}

\subsection{Mô tả đối tượng công nghệ}

Hệ thống cấp dầu nhiên liệu khởi động thường bao gồm các khâu chính: tiếp nhận dầu, lưu trữ trong bể, bơm chuyển, gia nhiệt, tuần hoàn hồi lưu và phân phối đến các đầu đốt. Dầu sau khi tiếp nhận từ nguồn ngoài được đưa vào bể đệm, sau đó bơm chuyển dầu đến các bể chứa chính. Tại đây, mức dầu được giám sát liên tục nhằm bảo đảm nguồn cấp ổn định cho hệ thống bơm hút.

Từ các bể chứa chính, dầu được bơm qua tuyến công nghệ chính đến bộ gia nhiệt. Mục đích của gia nhiệt là giảm độ nhớt động học của dầu xuống mức phù hợp với điều kiện vận chuyển qua đường ống và phun tơi tại đầu đốt. Đối với dầu HFO hoặc mazut, nhiệt độ làm việc thường phải được duy trì trong vùng đủ cao để bảo đảm tính chảy và tính phun sương, đồng thời không vượt quá giới hạn gây phân hủy nhiệt hoặc tạo cặn bám.

Sau bộ gia nhiệt, dầu nóng được dẫn đến manifold phân phối để chia đến các nhóm đầu đốt. Lưu lượng trên từng nhánh được điều tiết bởi van điều khiển lưu lượng. Một phần dòng dầu còn được hồi lưu về bể chứa nhằm duy trì tuần hoàn, tránh lắng cặn và bảo đảm dầu luôn ở trạng thái nhiệt ổn định. Đường hồi lưu cũng có vai trò bảo vệ bơm khỏi làm việc trong điều kiện thiếu tải hoặc lưu lượng quá nhỏ.

Sơ đồ công nghệ tổng quan được trình bày trong Hình~\ref{fig:so_do_cong_nghe}. Sơ đồ này chỉ ra mối quan hệ giữa khâu lưu trữ, gia nhiệt, điều tiết và hồi lưu; đồng thời là cơ sở để xây dựng sơ đồ đo lường và cấu trúc điều khiển ở các phần sau.

\begin{figure}[htbp]
    \centering
    \begin{tikzpicture}[
        scale=0.95,
        every node/.style={transform shape},
        equipment/.style={draw, rectangle, minimum width=2.2cm, minimum height=1.0cm, align=center, thick, fill=white, font=\sffamily\small},
        pipe/.style={thick, draw=cyan!70!black},
        arrow/.style={-Latex, thick, draw=cyan!70!black}
    ]
        \node[equipment] (buffer) at (0, 5) {Bể chứa\đệm};
        \node[equipment, minimum width=1.8cm] (tank1) at (4, 5) {Tank 1};
        \node[equipment, minimum width=1.8cm] (tank2) at (8, 5) {Tank 2};
        \node[equipment] (heater1) at (4, 1.2) {Gia nhiệt\dầu};
        \node[equipment] (heater2) at (8, 1.2) {Gia nhiệt\dầu};
        \draw[pipe] (3.8, 7) -- (8.2, 7);
        \node[above, font=\sffamily\small] at (6, 7) {Dầu hồi lưu từ lò};
        \draw[pipe] (6, 7) -- (6, 5);
        \draw[pipe] (tank1.east) -- (tank2.west);
        \draw[arrow] (buffer.east) -- (tank1.west);
        \draw[arrow] (tank1.south) -- (4, 4.1) -- (4, 2.0) -- (heater1.north);
        \draw[arrow] (heater1.south) -- (4, -1.2);
        \draw[arrow] (tank2.south) -- (8, 4.1) -- (8, 2.0) -- (heater2.north);
        \draw[arrow] (heater2.south) -- (8, -1.2);
        \draw[pipe] (4, -1.2) -- (8, -1.2);
        \draw[pipe] (6, -1.2) -- (6, -2.2) -- (1.3, -2.2) -- (1.3, -1.2);
        \draw[pipe] (1.3, -2.2) -- (1.3, -0.8);
        \draw[pipe] (1.3, -0.8) -- (0.5, -0.8) -- (0.5, 0);
        \draw[pipe] (1.3, -0.8) -- (2.1, -0.8) -- (2.1, 0);
        \draw[pipe] (0.5, 0) -- (2.1, 0);
        \draw[pipe] (1.3, 0) -- (1.3, 1.0);
        \draw[thick, draw=cyan!70!black, {Latex}-{Latex}] (0.0, 1.0) -- (2.6, 1.0);
        \node[above, font=\sffamily\small] at (1.3, 1.0) {Đến đầu đốt};
    \end{tikzpicture}
    \caption{Sơ đồ công nghệ tổng quan hệ thống cấp dầu nhiên liệu}
    \label{fig:so_do_cong_nghe}
\end{figure}

\subsection{Các thông số điều khiển và giám sát}

Trong hệ thống cấp dầu nhiên liệu, các thông số cần được theo dõi và điều khiển liên tục gồm lưu lượng, áp suất, nhiệt độ và mức dầu. Lưu lượng là biến công nghệ trọng tâm, phản ánh trực tiếp lượng nhiên liệu được đưa tới đầu đốt. Ở chế độ vận hành ổn định, lưu lượng phải bám theo giá trị đặt của hệ thống điều phối để bảo đảm cường độ cháy phù hợp với yêu cầu phụ tải.

Áp suất đường ống phân phối có vai trò duy trì điều kiện vận chuyển và phun tơi ổn định. Nếu áp suất quá thấp, lưu lượng tại đầu đốt dao động mạnh và quá trình phun sương suy giảm; nếu áp suất quá cao, nguy cơ rò rỉ và hư hỏng thiết bị tăng lên. Vì vậy, áp suất thường được giữ trong một dải làm việc hẹp để tối ưu hóa tính ổn định động lực học của hệ thống.

Nhiệt độ dầu ảnh hưởng trực tiếp đến độ nhớt, khả năng bơm và chất lượng phun. Đối với dầu nặng, khi nhiệt độ thấp, độ nhớt tăng mạnh khiến tổn thất thủy lực tăng và quá trình phun bị xấu đi. Khi nhiệt độ quá cao, dầu có thể lão hóa nhiệt, tạo cặn và ảnh hưởng tới độ bền thiết bị. Do đó, nhiệt độ là thông số vừa mang tính công nghệ vừa mang tính bảo vệ đối với toàn bộ tuyến dầu.

Mức dầu trong bể chứa là biến giám sát quan trọng để bảo đảm bơm làm việc trong điều kiện hút ổn định. Mức quá thấp có thể gây xâm thực và làm mất áp hút, còn mức quá cao làm tăng nguy cơ tràn và mất an toàn cháy nổ. Các tín hiệu mức vì vậy thường gắn với cơ chế cảnh báo và liên động bảo vệ.

\subsection{Vai trò của hệ thống và các yếu tố ảnh hưởng}

Hệ thống điều khiển lưu lượng dầu khởi động có ba vai trò chính. Thứ nhất, nó bảo đảm cung cấp nhiên liệu đúng theo nhu cầu của quá trình cháy, giúp lò hơi khởi động tin cậy và duy trì ngọn lửa ổn định trong các chế độ tải khác nhau. Thứ hai, hệ thống góp phần nâng cao hiệu suất cháy bằng cách duy trì lưu lượng cấp và áp suất phân phối ở vùng tối ưu, từ đó hỗ trợ quá trình hòa trộn nhiên liệu với không khí diễn ra hiệu quả hơn. Thứ ba, hệ thống là một phần của cơ chế bảo vệ an toàn, vì mọi sai lệch lớn về lưu lượng, áp suất hoặc mức dầu đều có thể dẫn tới mất ổn định cháy hoặc sự cố cháy nổ.

Chất lượng vận hành của hệ thống chịu ảnh hưởng bởi nhiều nguồn nhiễu. Trước hết là biến thiên độ nhớt của dầu do thay đổi nhiệt độ môi trường hoặc điều kiện gia nhiệt. Tiếp theo là dao động áp suất đầu ra của bơm và tổn thất trên tuyến ống, làm đặc tính lưu lượng thay đổi theo thời gian. Ngoài ra, sự mài mòn của van điều khiển và cặn bẩn trong dầu có thể làm thay đổi đặc tính tuyến tính hóa của cơ cấu chấp hành. Đây là những nguyên nhân dẫn đến sai lệch động học cần được xem xét khi mô hình hóa và thiết kế điều khiển.

\subsection{Thiết bị công nghệ chính}

Hệ thống được cấu thành từ các thiết bị đo lường, chấp hành và điều khiển. Thiết bị đo lưu lượng đóng vai trò cung cấp tín hiệu phản hồi chính cho vòng điều khiển. Trong các ứng dụng dầu nặng, cảm biến lưu lượng kiểu Coriolis hoặc oval gear thường được ưu tiên nhờ khả năng đo ổn định trong môi trường có độ nhớt biến thiên. Van điều khiển lưu lượng là cơ cấu chấp hành trực tiếp, có nhiệm vụ thay đổi độ mở để điều tiết lưu lượng cấp đến đầu đốt. Bộ điều khiển trung tâm, thường là PLC hoặc DCS, thực hiện xử lý tín hiệu, so sánh sai lệch và phát lệnh điều chỉnh cho van. Bên cạnh đó, các thiết bị phụ trợ như cảm biến áp suất, cảm biến nhiệt độ, công tắc mức và hệ thống báo động liên động bảo đảm hệ thống vận hành an toàn trong toàn bộ dải công nghệ.

\subsection{Sơ đồ đo lường, điều khiển và giám sát}

Từ đặc điểm công nghệ nêu trên, sơ đồ đo lường và điều khiển của hệ thống có thể tổ chức theo cấu trúc vòng kín cơ bản. Cảm biến lưu lượng FT đo lưu lượng thực tế trên tuyến cấp dầu và gửi tín hiệu về bộ điều khiển FIC. Bộ điều khiển so sánh giá trị đo với giá trị đặt, sau đó phát tín hiệu điều khiển đến van FCV để điều chỉnh độ mở. Khi lưu lượng thay đổi, quá trình cháy trong lò cũng biến đổi tương ứng; vì vậy, hệ thống giám sát trung tâm phải liên kết thông số lưu lượng với áp suất, nhiệt độ và mức dầu để bảo đảm điều kiện vận hành đồng bộ.

Bên cạnh vòng lưu lượng chính, các vòng phụ như điều khiển áp suất và nhiệt độ dầu có vai trò hỗ trợ làm ổn định trạng thái vật lý của lưu chất. Các tín hiệu báo mức thấp thấp, mức cao cao, áp suất cao hoặc nhiệt độ vượt ngưỡng được đưa vào logic bảo vệ liên động. Cấu trúc đo lường và điều khiển như vậy vừa đáp ứng yêu cầu điều hành vận hành, vừa tạo cơ sở rõ ràng cho việc mô hình hóa và tổng hợp bộ điều khiển ở các chương sau.

\subsection{Kết luận chương}

Chương này đã làm rõ cấu trúc công nghệ, các thông số vận hành chủ yếu và vai trò kỹ thuật của hệ thống cấp dầu nhiên liệu khởi động cho lò than phun. Qua phân tích có thể thấy đối tượng là một hệ quá trình có quán tính lớn, chịu ảnh hưởng mạnh của nhiệt độ, áp suất và đặc tính lưu chất, đồng thời đòi hỏi độ tin cậy cao về an toàn. Những đặc điểm đó là cơ sở để lựa chọn thiết bị đo, cơ cấu chấp hành và triết lý điều khiển phù hợp trong các chương tiếp theo.

\end{document}